\documentclass[a4paper]{article}
\usepackage[fontset=fandol]{ctex}
\usepackage{amsmath, amssymb}
\usepackage{graphicx}
\usepackage[margin=2.5cm]{geometry}
\usepackage[most]{tcolorbox}
\usepackage{etoolbox}
\usepackage{listings}
\usepackage{hyperref}
\usepackage{booktabs}
\usepackage{subcaption}
\usepackage{float}
\usepackage{tikz}
\IfFileExists{pgfplots.sty}{\usepackage{pgfplots}\pgfplotsset{compat=1.18}}{}
\newtcolorbox{knowledgebox}[1]{enhanced, colback=blue!5!white, colframe=blue!75!black, colbacktitle=blue!75!black, coltitle=white, fonttitle=\bfseries, title=#1, attach boxed title to top left={yshift=-2mm, xshift=2mm}, boxrule=1pt, sharp corners}
\newtcolorbox{importantbox}[1]{enhanced, colback=yellow!10!white, colframe=yellow!80!black, colbacktitle=yellow!80!black, coltitle=black, fonttitle=\bfseries, title=#1, sharp corners}
\newtcolorbox{warningbox}[1]{enhanced, colback=red!5!white, colframe=red!75!black, colbacktitle=red!75!black, coltitle=white, fonttitle=\bfseries, title=#1, sharp corners}
\lstset{language=Python,basicstyle=\ttfamily\small,keywordstyle=\color{blue},stringstyle=\color{red!60!black},commentstyle=\color{green!60!black},breaklines=true,frame=single,numbers=left,numberstyle=\tiny\color{gray},captionpos=b,extendedchars=false}
\newcommand{\notetitle}{End-to-End ASR: RNN-T}
\newcommand{\noteauthors}{Codex 基于课程页、公开视频字幕、coverage ledger 与补充官方材料整理}
\newcommand{\notedate}{\today}
\newcommand{\videochannel}{WAVLab}
\newcommand{\videopublishdate}{2023-11-13}
\newcommand{\videoduration}{unknown}
\newcommand{\videourl}{https://www.youtube.com/watch?v=BQBOu9BOFpc}
\newcommand{\videocoverpath}{cover.jpg}
\begin{document}
\begin{titlepage}
\centering
{\Large 课程笔记\par}
\vspace{1.2cm}
{\huge\bfseries \notetitle\par}
\vspace{0.8cm}
{\large \noteauthors\par}
\vspace{0.3cm}
{\large \notedate\par}
\vspace{1.2cm}
\ifdefempty{\videocoverpath}{{\small 公开视频封面不可用。\par}}{\includegraphics[width=0.82\textwidth,height=0.45\textheight,keepaspectratio]{\videocoverpath}\par}
\vfill
\begin{tcolorbox}[width=0.9\textwidth, colback=black!2!white, colframe=black!60, sharp corners]
\textbf{视频作者/频道}：\videochannel\par
\textbf{发布日期}：\videopublishdate\par
\textbf{视频时长}：\videoduration\par
\textbf{视频链接}：\href{\videourl}{\nolinkurl{\videourl}}
\end{tcolorbox}
\end{titlepage}
\tableofcontents
\newpage
\section{RNN-T overview}
根据课程页主题、公开视频字幕以及本讲的 coverage ledger，RNN-T overview 这一部分主要围绕 End-to-End ASR: RNN-T | End-to-End ASR: RNN-T；RNN-T 的总体结构、它和 CTC / attention 的关系，以及为什么它适合 end-to-end ASR。 展开。
写作上保持 coverage-first 原则：先把老师在这一阶段真正讲过的问题设定、模型部件、训练或推理逻辑讲清，再压缩成适合复习的结构。

\begin{figure}[H]
\centering
\includegraphics[width=0.92\textwidth]{figures/coverage_map.png}
\caption{本讲 coverage map：按 coverage\_units.jsonl 归纳出的核心主题与章节映射。}
\end{figure}

\begin{figure}[H]
\centering
\includegraphics[width=0.92\textwidth]{figures/timeline.png}
\caption{本讲 timeline：根据字幕时间范围归纳的主题推进顺序。}
\end{figure}

\subsection{End-to-End ASR: RNN-T}
从证据位置看，这一部分对应的课程证据区间是 2023-11-13。课程在这里强调的核心内容可以概括为：End-to-End ASR: RNN-T | End-to-End ASR: RNN-T。
课程把 RNN-T 放在 streaming ASR 的语境里讨论，强调它保留了单调约束，却让模型可以更灵活地联合时间和输出历史。

在本章结构中，这一单元被并入“RNN-T overview”，目的是把同一阶段的定义、机制和权衡集中讲清。
\[
P(Y \mid X) = \sum_{\pi \in \mathcal{A}(X,Y)} \prod_{(t,u)\in\pi} P(k_{t,u} \mid h_t, g_u)
\]

符号说明如下：
\begin{itemize}
\item $X$：输入语音特征序列；
\item $Y$：目标标签序列；
\item $\mathcal{A}(X,Y)$：满足转导约束的对齐集合；
\item $h_t$：encoder 在时间步 $t$ 的表示；
\item $g_u$：prediction network 在输出步 $u$ 的表示。
\end{itemize}

\subsection{RNN-T 的总体结构、它和 CTC / attention 的关系，以及为什么它适合 end-to-end ASR}
从证据位置看，这一部分对应的课程证据区间是 00:00:00,880--00:07:57,479。课程在这里强调的核心内容可以概括为：RNN-T 的总体结构、它和 CTC / attention 的关系，以及为什么它适合 end-to-end ASR。。
CTC 的关键在于保持单调对齐假设，同时把对齐路径视作潜变量并在训练时进行求和。

在本章结构中，这一单元被并入“RNN-T overview”，目的是把同一阶段的定义、机制和权衡集中讲清。
\[
P(Y \mid X) = \sum_{\pi \in \mathcal{B}^{-1}(Y)} \prod_{t=1}^{T} P(\pi_t \mid X)
\]

符号说明如下：
\begin{itemize}
\item $X$：输入语音特征序列；
\item $Y$：输出标签序列；
\item $\pi$：带有 blank 的对齐路径；
\item $\mathcal{B}$：把路径压缩为最终标签序列的映射。
\end{itemize}

\subsection{本章小结}
RNN-T overview 这一部分把 End-to-End ASR: RNN-T 中相邻的教学单元串成了一条连续叙事：既交代了这一阶段要解决的问题，也交代了方法、推理或工程权衡为什么要这样组织。

\section{Encoder-predictor-joiner architecture}
根据课程页主题、公开视频字幕以及本讲的 coverage ledger，Encoder-predictor-joiner architecture 这一部分主要围绕 encoder / prediction network / joiner 三部分如何共同定义 transducer 计算。 展开。
写作上保持 coverage-first 原则：先把老师在这一阶段真正讲过的问题设定、模型部件、训练或推理逻辑讲清，再压缩成适合复习的结构。

\subsection{encoder / prediction network / joiner 三部分如何共同定义 transducer 计算}
从证据位置看，这一部分对应的课程证据区间是 00:07:55,639--00:16:23,759。课程在这里强调的核心内容可以概括为：encoder / prediction network / joiner 三部分如何共同定义 transducer 计算。。
课程把 RNN-T 放在 streaming ASR 的语境里讨论，强调它保留了单调约束，却让模型可以更灵活地联合时间和输出历史。

在本章结构中，这一单元被并入“Encoder-predictor-joiner architecture”，目的是把同一阶段的定义、机制和权衡集中讲清。
\[
P(Y \mid X) = \sum_{\pi \in \mathcal{A}(X,Y)} \prod_{(t,u)\in\pi} P(k_{t,u} \mid h_t, g_u)
\]

符号说明如下：
\begin{itemize}
\item $X$：输入语音特征序列；
\item $Y$：目标标签序列；
\item $\mathcal{A}(X,Y)$：满足转导约束的对齐集合；
\item $h_t$：encoder 在时间步 $t$ 的表示；
\item $g_u$：prediction network 在输出步 $u$ 的表示。
\end{itemize}

\subsection{本章小结}
Encoder-predictor-joiner architecture 这一部分把 End-to-End ASR: RNN-T 中相邻的教学单元串成了一条连续叙事：既交代了这一阶段要解决的问题，也交代了方法、推理或工程权衡为什么要这样组织。

\section{Alignment factorization}
根据课程页主题、公开视频字幕以及本讲的 coverage ledger，Alignment factorization 这一部分主要围绕 从 acoustic model 到 sequence-level objective，说明 RNN-T 的对齐分解。 展开。
写作上保持 coverage-first 原则：先把老师在这一阶段真正讲过的问题设定、模型部件、训练或推理逻辑讲清，再压缩成适合复习的结构。

\subsection{从 acoustic model 到 sequence-level objective，说明 RNN-T 的对齐分解}
从证据位置看，这一部分对应的课程证据区间是 00:16:20,720--00:23:46,039。课程在这里强调的核心内容可以概括为：从 acoustic model 到 sequence-level objective，说明 RNN-T 的对齐分解。。
课程把 RNN-T 放在 streaming ASR 的语境里讨论，强调它保留了单调约束，却让模型可以更灵活地联合时间和输出历史。

在本章结构中，这一单元被并入“Alignment factorization”，目的是把同一阶段的定义、机制和权衡集中讲清。
\[
P(Y \mid X) = \sum_{\pi \in \mathcal{A}(X,Y)} \prod_{(t,u)\in\pi} P(k_{t,u} \mid h_t, g_u)
\]

符号说明如下：
\begin{itemize}
\item $X$：输入语音特征序列；
\item $Y$：目标标签序列；
\item $\mathcal{A}(X,Y)$：满足转导约束的对齐集合；
\item $h_t$：encoder 在时间步 $t$ 的表示；
\item $g_u$：prediction network 在输出步 $u$ 的表示。
\end{itemize}

\subsection{本章小结}
Alignment factorization 这一部分把 End-to-End ASR: RNN-T 中相邻的教学单元串成了一条连续叙事：既交代了这一阶段要解决的问题，也交代了方法、推理或工程权衡为什么要这样组织。

\section{Training and approximation}
根据课程页主题、公开视频字幕以及本讲的 coverage ledger，Training and approximation 这一部分主要围绕 teacher forcing、近似训练与优化细节如何影响 transducer 学习。 展开。
写作上保持 coverage-first 原则：先把老师在这一阶段真正讲过的问题设定、模型部件、训练或推理逻辑讲清，再压缩成适合复习的结构。

\subsection{teacher forcing、近似训练与优化细节如何影响 transducer 学习}
从证据位置看，这一部分对应的课程证据区间是 00:23:41,440--00:31:40,799。课程在这里强调的核心内容可以概括为：teacher forcing、近似训练与优化细节如何影响 transducer 学习。。
课程把 RNN-T 放在 streaming ASR 的语境里讨论，强调它保留了单调约束，却让模型可以更灵活地联合时间和输出历史。

在本章结构中，这一单元被并入“Training and approximation”，目的是把同一阶段的定义、机制和权衡集中讲清。
\[
P(Y \mid X) = \sum_{\pi \in \mathcal{A}(X,Y)} \prod_{(t,u)\in\pi} P(k_{t,u} \mid h_t, g_u)
\]

符号说明如下：
\begin{itemize}
\item $X$：输入语音特征序列；
\item $Y$：目标标签序列；
\item $\mathcal{A}(X,Y)$：满足转导约束的对齐集合；
\item $h_t$：encoder 在时间步 $t$ 的表示；
\item $g_u$：prediction network 在输出步 $u$ 的表示。
\end{itemize}

\subsection{本章小结}
Training and approximation 这一部分把 End-to-End ASR: RNN-T 中相邻的教学单元串成了一条连续叙事：既交代了这一阶段要解决的问题，也交代了方法、推理或工程权衡为什么要这样组织。

\section{Decoding and streaming}
根据课程页主题、公开视频字幕以及本讲的 coverage ledger，Decoding and streaming 这一部分主要围绕 beam search decoding、language model integration 与在线 / streaming 场景。 展开。
写作上保持 coverage-first 原则：先把老师在这一阶段真正讲过的问题设定、模型部件、训练或推理逻辑讲清，再压缩成适合复习的结构。

\subsection{beam search decoding、language model integration 与在线 / streaming 场景}
从证据位置看，这一部分对应的课程证据区间是 00:31:37,919--00:39:49,560。课程在这里强调的核心内容可以概括为：beam search decoding、language model integration 与在线 / streaming 场景。。
Search 这一讲把前面所有模型真正接到了 decoding pipeline 上，也就是把训练得到的分数转成最终转写结果。

在本章结构中，这一单元被并入“Decoding and streaming”，目的是把同一阶段的定义、机制和权衡集中讲清。

\subsection{本章小结}
Decoding and streaming 这一部分把 End-to-End ASR: RNN-T 中相邻的教学单元串成了一条连续叙事：既交代了这一阶段要解决的问题，也交代了方法、推理或工程权衡为什么要这样组织。

\section{Tradeoffs and recap}
根据课程页主题、公开视频字幕以及本讲的 coverage ledger，Tradeoffs and recap 这一部分主要围绕 RNN-T 的优势、局限与和前述 attention / CTC 方法的对比收束。 展开。
写作上保持 coverage-first 原则：先把老师在这一阶段真正讲过的问题设定、模型部件、训练或推理逻辑讲清，再压缩成适合复习的结构。

\subsection{RNN-T 的优势、局限与和前述 attention / CTC 方法的对比收束}
从证据位置看，这一部分对应的课程证据区间是 00:39:46,880--00:46:44,839。课程在这里强调的核心内容可以概括为：RNN-T 的优势、局限与和前述 attention / CTC 方法的对比收束。。
CTC 的关键在于保持单调对齐假设，同时把对齐路径视作潜变量并在训练时进行求和。

在本章结构中，这一单元被并入“Tradeoffs and recap”，目的是把同一阶段的定义、机制和权衡集中讲清。
\[
P(Y \mid X) = \sum_{\pi \in \mathcal{B}^{-1}(Y)} \prod_{t=1}^{T} P(\pi_t \mid X)
\]

符号说明如下：
\begin{itemize}
\item $X$：输入语音特征序列；
\item $Y$：输出标签序列；
\item $\pi$：带有 blank 的对齐路径；
\item $\mathcal{B}$：把路径压缩为最终标签序列的映射。
\end{itemize}

\subsection{本章小结}
Tradeoffs and recap 这一部分把 End-to-End ASR: RNN-T 中相邻的教学单元串成了一条连续叙事：既交代了这一阶段要解决的问题，也交代了方法、推理或工程权衡为什么要这样组织。

\section{总结与延伸}
End-to-End ASR: RNN-T 这一讲在整门《Speech Recognition and Understanding》课程中的位置，是把 E2E ASR: RNN-T 讲成一个可复用的建模模块。
从教材化角度看，本讲最重要的不是记住单一术语，而是理解它如何嵌入完整的 automatic speech recognition pipeline：输入语音如何被表示、模型如何被训练、解码如何得到最终转写，以及这些设计分别带来什么假设和权衡。
由于公开源以课程页和公开视频字幕为主、缺少官方 slide PDF，本章在正文里尽量把术语、方法和工程语境解释完整，并把缺失项留在 omission 和 manifest 里，而不是凭空补写。
\end{document}
